\section{Actual diagonal lifts and exceptional direction normalization}
\label{sec:hironaka-normal-crossings-diagonal-lift-orbits}

Fix the actual coordinate centre $L_k=\{z_i=0:i\geq k\}$ in the
lifted coordinate model and an actual blowup $\sigma:Z\to\mathbb C^n$
along $L_k$. All monomial-chart lifts below are the genuine unique
holomorphic lifts supplied by this blowup's universal property. We
prove geometric point identities, without asserting a finite whole
canonical value image or replacing the original one-step history.

\begin{lemma}[Diagonal biholomorphisms and the whole monomial equation]
\label{hir:general-nc-diagonal-monomial-chart-equation}
\leanok
\lean{Hironaka.HironakaNormalCrossingsDiagonalLiftOrbits.diagonalEquiv,Hironaka.HironakaNormalCrossingsDiagonalLiftOrbits.diagonal_preimage_centre,Hironaka.HironakaNormalCrossingsDiagonalLiftOrbits.diagonal_preimage_carrier,Hironaka.HironakaNormalCrossingsDiagonalLiftOrbits.monomial_chartScale}
Let every $a_i$ be nonzero. The diagonal map $D_a(z)_i=a_i z_i$ is an
actual biholomorphism preserving $L_k$ and every selected coordinate
hyperplane union. In the monomial chart of pivot $j\geq k$, put
$T_{a,j}(w)_j=a_jw_j$, $T_{a,j}(w)_i=(a_i/a_j)w_i$ for $i\geq k$, $i\ne j$,
and $T_{a,j}(w)_i=a_iw_i$ for $i<k$. Then the whole polynomial maps satisfy
$m_j\circ T_{a,j}=D_a\circ m_j$ at every chart point.
\end{lemma}
\begin{proof}
The diagonal inverse has coefficients $a_i^{-1}$, and coordinatewise
linearity makes both maps holomorphic. Nonzero multiplication preserves
each coordinate zero predicate, so both the whole centre and every
coordinate union are preserved. For a nonpivot normal coordinate, the
left side equals $(a_i/a_j)w_i\,a_jw_j=a_iw_iw_j$, using $a_j\ne0$.
At the pivot and at every free coordinate, the monomial map retains
that coordinate and both sides equal $a_iw_i$. These coordinatewise
equalities prove the entire map equation, including all zero coordinates.
\end{proof}



\begin{theorem}[The actual diagonal lift has the actual chart point identity]
\label{hir:general-nc-actual-diagonal-lift-chart-identity}
\leanok
\lean{Hironaka.HironakaNormalCrossingsDiagonalLiftOrbits.exists_diagonal_history_lift_chart}
The actual diagonal base biholomorphism has an actual whole one-step
history lift $E_a$ on $Z$. For its genuine monomial lift $\ell_j$,
$E_a(\ell_j(w))=\ell_j(T_{a,j}(w))$ at every actual chart point.
\uses{hir:general-nc-one-step-history-universal-comparison,hir:general-nc-diagonal-monomial-chart-equation}
\end{theorem}
\begin{proof}
Whole centre preservation and the actual universal comparison of the
two genuine blowups construct an actual lifted biholomorphism and its
whole one-step history isomorphism. The exact projection equation gives
$\sigma\circ E_a\circ\ell_j=D_a\circ m_j$. The preceding whole monomial
equation also gives $\sigma\circ\ell_j\circ T_{a,j}=D_a\circ m_j$.
Both compared maps are actual holomorphic maps on the entire chart.
The pullback of the centre ideal along $m_j$ is invertible, generated
by its pivot coordinate. Composing the base with a centre-preserving
diagonal biholomorphism retains that invertibility. Thus uniqueness in
the actual blowup universal property applies to these two maps with the
same projection and makes them equal. Evaluating at $w$ proves the
asserted actual point identity. Equality of projections alone is not
used as a substitute for this uniqueness argument.
\end{proof}

\begin{theorem}[Actual zero--one exceptional direction representatives]
\label{hir:general-nc-actual-binary-direction-normalization}
\leanok
\lean{Hironaka.HironakaNormalCrossingsDiagonalLiftOrbits.exists_binary_direction_history_lift}
Every actual point over zero can be carried by an actual diagonal
one-step history isomorphism to a genuine monomial-chart point whose
free and pivot coordinates are zero and whose remaining normal
coordinates belong to $\{0,1\}$. The pivot and its nonzero support are
actual finite coordinate data.
\uses{hir:general-nc-actual-diagonal-lift-chart-identity}
\end{theorem}
\begin{proof}
Actual monomial-chart surjectivity gives a pivot $j\geq k$ and a genuine
chart point $w$ mapping to the specified blown-up point. Its exact
projection is zero, so the pivot coordinate and every free coordinate
of $w$ are zero. Put $a_j=1$. At every other normal coordinate, put
$a_i=1$ when $w_i=0$ and $a_i=w_i^{-1}$ otherwise; put $a_i=1$ at free
coordinates. Every chosen coefficient is nonzero. The resulting
$T_{a,j}(w)$ has precisely the asserted zero--one coordinates. The
preceding actual lifted point identity sends the original blown-up
point to the genuine chart lift of this exact finite-support vector.
The proof constructs the base coefficients, actual lifted map and
actual point equality internally. No invariant value, orbit
equivalence, history isomorphism or finite-value certificate is input.
To deduce finitely many whole canonical values, a separate comparison
with a common literal canonical model history is still required.
\end{proof}

\begin{corollary}[An actual finite set of geometric fibre representatives]
\label{hir:general-nc-finite-geometric-fibre-orbit-representatives}
\leanok
\lean{Hironaka.HironakaNormalCrossingsDiagonalLiftOrbits.exists_finite_geometric_orbit_representatives}
There is an actual finite set $Q\subset Z$, depending only on the fixed
actual coordinate blowup and its genuine chart lifts, such that each
actual point over zero is carried into $Q$ by a genuine diagonal
one-step history isomorphism. This is a geometric statement; no finite
whole invariant value image is asserted.
\uses{hir:general-nc-actual-binary-direction-normalization}
\end{corollary}
\begin{proof}
For every pivot $j\geq k$ and every subset $S$ of the finite coordinate
index type, take the actual chart lift of its binary support vector.
Let $Q$ be the image of this finite indexing family. It is finite even
when distinct indices give the same actual point. The preceding theorem
constructs, for each point over zero, an actual pivot and support and
a genuine diagonal history isomorphism sending it to precisely that
indexed chart point. Hence its image belongs to this same fixed $Q$.
If the indexing family is empty, actual monomial-chart surjectivity
implies there is no point over zero, so the quantified conclusion
still follows from that preceding theorem.
\end{proof}
